\documentclass[11pt]{amsart}

\usepackage[margin=1.05in]{geometry}
\usepackage{amsmath,amssymb,amsthm,mathtools}
\usepackage{enumitem}
\usepackage[T1]{fontenc}
\usepackage{lmodern}
\usepackage{hyperref}

\newtheorem{theorem}{Theorem}[section]
\newtheorem{proposition}[theorem]{Proposition}
\newtheorem{lemma}[theorem]{Lemma}
\newtheorem{corollary}[theorem]{Corollary}
\newtheorem{definition}[theorem]{Definition}
\newtheorem{opentheorem}[theorem]{Open theorem}
\newtheorem{target}[theorem]{Certificate target}
\newtheorem{question}[theorem]{Question}
\newtheorem{remark}[theorem]{Remark}

\newcommand{\kbar}{\overline{\mathbb F}_5}
\newcommand{\PP}{\mathbb P}
\newcommand{\Div}{\operatorname{div}}
\newcommand{\Norm}{\operatorname{Norm}}
\newcommand{\Tr}{\operatorname{Tr}}
\newcommand{\Pic}{\operatorname{Pic}}
\newcommand{\Spec}{\operatorname{Spec}}
\newcommand{\Aut}{\operatorname{Aut}}
\newcommand{\ord}{\operatorname{ord}}
\newcommand{\Comm}{\operatorname{Comm}}
\newcommand{\cO}{\mathcal O}
\newcommand{\F}{\mathbb F}

\title{A conditional reduction of a characteristic-five common-cover obstruction}
\author{June 2026}

\begin{document}
\maketitle

\begin{abstract}
Let \(k=\overline{\mathbb F}_5\) and \(S=\PP^1_k(31,31,31)\).  The common-cover
counterexample considered here is reduced to the visibility of finite etale
self-correspondences of \(S\).  The reduction is then carried to the curve
level: a non-visible correspondence is assumed to yield a particular
genus-\(11\), degree-\(35\) pair of three-point maps, and the remaining work is
to exclude three explicit branches of this pair.  One branch is further
reduced to two stated computer-algebra claims.  These four open theorems
together imply the visibility theorem for \(S\), hence the desired
common-cover obstruction.
\end{abstract}

\tableofcontents

\section{Statement of the Reduction}

The original question is the following.

\begin{question}[Common-cover question over \(\kbar\)]
Do any two smooth proper connected curves of genus at least \(2\) over
\(k=\kbar\) admit a common connected finite etale cover?
\end{question}

The proposed counterexample is \(Y:y^{31}=x(x-1)\).  Put
\(S=\PP^1_k(31,31,31)\).  Section~\ref{sec:global-route} proves that it is
enough to show every connected finite etale self-correspondence of \(S\) is
visible, in the sense of Definition~\ref{def:visible}.  The curve-level part
of the note reduces this visibility theorem to Open
Theorems~\ref{open:profile-extraction}, \ref{open:entry-one},
\ref{open:no-highpoint}, and \ref{open:double-fiber}; the formal implication
is Theorem~\ref{thm:four-open-imply-comm}.

\section{The global stack route}
\label{sec:global-route}

Let \(k=\kbar\).  The stack
\[
  S=\PP^1_k(31,31,31)
\]
is the smooth proper stacky curve with coarse space \(\PP^1_k\) and stabilizer
order \(31\) at \(0,1,\infty\).

\begin{lemma}[The cyclic atlas]
\label{lem:atlas}
Let
\[
  Y:\quad y^{31}=x(x-1)
\]
and let \(C_{31}\) act by \(y\mapsto \zeta y\).  Then
\[
  S=[Y/C_{31}],
\]
the map \(Y\to S\) is a representable finite etale \(C_{31}\)-torsor, and
\(g(Y)=15\).
\end{lemma}

\begin{proof}
The coarse map \(Y\to\PP^1_x\) has degree \(31\) and is totally ramified over
\(0,1,\infty\).  Riemann--Hurwitz gives
\[
  2g(Y)-2=31(-2)+3(31-1)=28,
\]
so \(g(Y)=15\).  The quotient stack \([Y/C_{31}]\) remembers the three coarse
ramification stabilizers; therefore the quotient map \(Y\to [Y/C_{31}]\) is
finite etale in the stack sense and representable.
\end{proof}

The three stacky points have the same order, so \(S_3\) acts on \(S\) by
permuting them.  The quotient is
\[
  S_0=S/S_3\simeq \PP^1_k(2,3,62),
\]
and the quotient map is denoted
\[
  \pi_0:S\longrightarrow S_0.
\]
It is finite etale of degree \(6\) as a morphism of stacky curves: the branch
orders of the coarse \(S_3\)-quotient are exactly absorbed by the stack
indices \(2,3,62\) on the target.

\begin{definition}[Finite etale self-correspondence]
A finite etale self-correspondence of \(S\) is a connected smooth proper
Deligne--Mumford curve \(T\) together with two representable finite etale maps
\[
  u,v:T\to S.
\]
\end{definition}

\begin{definition}[Visible self-correspondence]
\label{def:visible}
A finite etale self-correspondence \(u,v:T\to S\) is visible if
\[
  \pi_0\circ u\simeq \pi_0\circ v.
\]
Equivalently, \((u,v):T\to S\times S\) factors through the fiber product
\[
  S\times_{S_0}S.
\]
Since \(S\to S_0\) is the \(S_3\)-quotient and \(T\) is connected, this is the
same as saying that \(v=\sigma\circ u\) for one fixed \(\sigma\in S_3\).
\end{definition}

The assertion that all finite etale self-correspondences of \(S\) are visible
is the visibility theorem for \(S\).  In the usual shorthand it is written
\[
  \Comm_{\mathrm{alg},k}(S)=\Aut(S)=S_3.
\]

\begin{lemma}[No genus-three etale cover of \(S_0\)]
\label{lem:no-genus3-S0}
No smooth proper connected genus-three curve admits a finite etale map to
\[
  S_0=\PP^1_k(2,3,62).
\]
\end{lemma}

\begin{proof}
For a stacky projective line \(\PP^1(a,b,c)\),
\[
  \deg K_{\PP^1(a,b,c)}
  =
  -2+\left(1-\frac1a\right)+\left(1-\frac1b\right)
      +\left(1-\frac1c\right).
\]
For \((a,b,c)=(2,3,62)\) this is \(14/93\).  If \(X\to S_0\) were finite
etale of degree \(d\) and \(g(X)=3\), then
\[
  4=\deg K_X=d\cdot\frac{14}{93},
\]
so \(d=186/7\), impossible.
\end{proof}

\begin{proposition}[Conditional common-cover counterexample]
\label{prop:conditional-common-cover}
Assume every finite etale self-correspondence of \(S\) is visible.  Then
\(Y:y^{31}=x(x-1)\) has no connected common finite etale cover with any smooth
proper connected genus-three curve over \(k\).
\end{proposition}

\begin{proof}
Suppose, to the contrary, that \(Z\to X\) and \(Z\to Y\) are connected finite
etale covers with \(g(X)=3\).  Compose \(Z\to Y\) with the atlas torsor
\(Y\to S\), and call the resulting finite etale map \(b:Z\to S\).

Let
\[
  R=Z\times_X Z
\]
with projections \(p_1,p_2:R\rightrightarrows Z\).  For each connected
component \(R_\alpha\) of \(R\), the two maps
\[
  b\circ p_1,\ b\circ p_2:R_\alpha\to S
\]
form a finite etale self-correspondence of \(S\).  By
Definition~\ref{def:visible},
\[
  \pi_0\circ b\circ p_1=\pi_0\circ b\circ p_2
\]
on every \(R_\alpha\), hence on all of \(R\).  Thus
\(\pi_0\circ b:Z\to S_0\) is invariant under the finite etale equivalence
relation \(R\rightrightarrows Z\).  Since \(Z\to X\) is finite etale
surjective, effective descent for finite etale morphisms gives a finite etale
map
\[
  X\to S_0.
\]
This contradicts Lemma~\ref{lem:no-genus3-S0}.
\end{proof}

\begin{remark}
The proposition reduces the common-cover counterexample to the algebraic
visibility theorem for \(S\).  The rest of the note reduces that theorem to
the four open theorems assembled in
Theorem~\ref{thm:four-open-imply-comm}.
\end{remark}

\section{Profile Pairs}
\label{sec:profile-pairs}

\begin{definition}[Three-point profile map]
\label{def:one-profile-map}
Let \(n,m\ge 1\).  A separable function \(h:C\to\PP^1\) on a smooth proper
connected curve has three-point profile \((n;m)\) if \(\deg h=n+m\), all
ramification of \(h\) lies over \(0,1,\infty\), and for each
\(a\in\{0,1,\infty\}\) one has
\[
  h^{-1}(a)=nP_a(h)+E_a(h),
\]
where \(P_a(h)\) is one point and \(E_a(h)\) is a reduced effective divisor of
degree \(m\), disjoint from the two other \(E\)-divisors for the same map.
In the rest of the note, ``profile map'' means three-point profile
\((31;4)\).
\end{definition}

\begin{lemma}[Genus of a \((31;4)\)-profile map]
\label{lem:profile-genus}
If \(C\) admits a \((31;4)\)-profile map, then \(g(C)=11\).
\end{lemma}

\begin{proof}
The ramification contribution over each of \(0,1,\infty\) is \(31-1=30\).
Since there is no other ramification,
\[
  2g(C)-2=35(-2)+3\cdot 30=20,
\]
so \(g(C)=11\).
\end{proof}

\begin{definition}[Profile-\(4\) two-map pair]
\label{def:profile4-pair}
A profile-\(4\) two-map pair is a smooth proper connected curve \(C/k\) with
two \((31;4)\)-profile maps \(x,r:C\to\PP^1\), such that
\[
  k(C)=k(x,r),
\]
the twelve unramified boundary points agree as a reduced divisor
\[
  U=E_0+E_1+E_\infty=F_0+F_1+F_\infty,
\]
where \(E_i=E_i(x)\) and \(F_j=E_j(r)\), and the incidence matrix
\[
  M_{ij}:=\deg(E_i\cap F_j)
\]
is
\[
  M=
  \begin{pmatrix}
  0&1&3\\
  1&2&1\\
  3&1&0
  \end{pmatrix}
  \qquad (i,j\in\{0,1,\infty\}).
\]
\end{definition}

\begin{lemma}[Divisors attached to a profile-\(4\) pair]
\label{lem:profile-divisors}
Let \((C,x,r)\) be a profile-\(4\) pair.  With
\[
  P_i=P_i(x),\qquad Q_j=P_j(r),
\]
one has
\[
 \Div(x)=31P_0+E_0-31P_\infty-E_\infty.
\]
Moreover
\[
 \Div(x-1)=31P_1+E_1-31P_\infty-E_\infty,
\]
and
\[
 \Div(r)=31Q_0+F_0-31Q_\infty-F_\infty,\qquad
 \Div(r-1)=31Q_1+F_1-31Q_\infty-F_\infty.
\]
\end{lemma}

\begin{proof}
For \(x\), the zero divisor is the scheme-theoretic fiber over \(0\), namely
\[
  x^{-1}(0)=31P_0+E_0,
\]
and the pole divisor is the fiber over \(\infty\),
\[
  x^{-1}(\infty)=31P_\infty+E_\infty.
\]
Their difference is \(\Div(x)\).  The same argument applied to \(x-1\) uses
the fiber of \(x\) over \(1\) for the zeros and the same pole fiber as \(x\).
The two formulas for \(r\) are identical with \(P,E,x\) replaced by
\(Q,F,r\).
\end{proof}

Write the nonempty atoms of the incidence decomposition as
\[
\begin{array}{lll}
 A=U_{0,1}, & B=U_{0,\infty}, & C=U_{1,0},\\
 D=U_{1,1}, & E=U_{1,\infty}, & F=U_{\infty,0},\\
 G=U_{\infty,1}. &&
\end{array}
\]
Their degrees are
\[
  \deg A=\deg C=\deg E=\deg G=1,\qquad
  \deg D=2,\qquad
  \deg B=\deg F=3.
\]
The atoms \(U_{0,0}\) and \(U_{\infty,\infty}\) are empty.

\begin{opentheorem}[A. Profile extraction from a non-visible correspondence]
\label{open:profile-extraction}
Let
\[
  u,v:T\to S=\PP^1_k(31,31,31)
\]
be a connected representable finite etale self-correspondence such that
\[
  \pi_0\circ u\not\simeq \pi_0\circ v.
\]
Then, after possibly interchanging \(u\) and \(v\) and composing one of them
with a visible \(S_3\)-symmetry of \(S\), there is a smooth proper connected
curve \(C/k\) and separable functions \(x,r\in k(C)\) such that
\((C,x,r)\) is a profile-\(4\) two-map pair in the sense of
Definition~\ref{def:profile4-pair}.
\end{opentheorem}

\begin{remark}
Open Theorem~\ref{open:profile-extraction} is the global-to-curve reduction:
it replaces a non-visible stack correspondence by the explicit two-map
incidence datum of Definition~\ref{def:profile4-pair}.
\end{remark}

\begin{remark}[Why profile \(4\)]
Among \(3\times3\) nonnegative integer matrices with all row and column sums
\(4\), modulo independent relabeling of rows and columns, there are nine
incidence profiles.  The profile above is the smallest of the six profiles in
which no \(r\)-boundary block is equal to \(E_\infty\).  This matters for the
Kummer tests in Section~\ref{sec:kummer}.  The unproved content in
Open Theorem~\ref{open:profile-extraction} is that, after the visible
\(S_3\)-symmetries and the preceding reductions, no other incidence profile
has to be considered.
\end{remark}

\section{Logarithmic Differentials and Kummer Tests}
\label{sec:kummer}

We use the following logarithmic identities.

\subsection{Logarithmic forms}

Let \(h\) be one profile map.  For \(c\in\{0,1,\infty\}\), define
\[
  \Omega_c(h)=\frac{dh}{(h-a)(h-b)},\qquad \{a,b,c\}=\{0,1,\infty\}.
\]
Thus
\[
  \Omega_0(h)=\frac{dh}{h-1},\qquad
  \Omega_1(h)=\frac{dh}{h},\qquad
  \Omega_\infty(h)=\frac{dh}{h(h-1)}.
\]

\begin{lemma}[Divisor of \(\Omega_c\)]
\label{lem:omega-divisor}
For \(\{a,b,c\}=\{0,1,\infty\}\),
\[
  \Div\Omega_c(h)
  =30P_c(h)-P_a(h)-P_b(h)-E_a(h)-E_b(h).
\]
\end{lemma}

\begin{proof}
At \(P_c(h)\), use \(h-c\) as a target coordinate if \(c\ne\infty\), and use
\(1/h\) if \(c=\infty\).  In either coordinate the local form of the map is
\(t\mapsto t^{31}\), so \(dh\) contributes order \(30\) relative to a local
parameter on \(C\), while neither denominator factor vanishes.  Hence
\(\Omega_c(h)\) has order \(30\) at \(P_c(h)\).

At \(P_a(h)\), the corresponding denominator factor has order \(31\), while
\(dh\) has order \(30\), so the order is \(-1\); the same argument applies at
\(P_b(h)\).  At each point of \(E_a(h)\) or \(E_b(h)\), the map is unramified
and the corresponding denominator factor has a simple zero, so the order is
\(-1\).  Away from the three boundary fibers the denominator is nonzero and
the map is unramified, so there are no further zeros or poles.  The displayed
divisor also has the required degree
\[
  30-1-1-4-4=20=2g(C)-2.
\]
\end{proof}

At an unramified boundary point with boundary label \(\ell\), the residue of
\(\Omega_c(h)\) is
\[
\begin{array}{c|ccc}
       & \ell=0 & \ell=1 & \ell=\infty\\ \hline
c=0    & 0&1&-1\\
c=1    & 1&0&-1\\
c=\infty&-1&1&0 .
\end{array}
\]
For example, near a point with \(h=0\), the form \(dh/h\) has residue \(1\),
while \(dh/(h-1)\) is regular.

Multiplying the three logarithmic forms for one map gives
\[
  \Div\bigl(\Omega_0(h)\Omega_1(h)\Omega_\infty(h)\bigr)
  =28(P_0(h)+P_1(h)+P_\infty(h))-2U_h,
\]
where \(U_h=E_0(h)+E_1(h)+E_\infty(h)\).  Hence a profile-\(4\) pair satisfies
the Picard constraint
\[
  28\bigl((P_0+P_1+P_\infty)-(Q_0+Q_1+Q_\infty)\bigr)=0
  \quad\text{in }\Pic^0(C).
\]

\subsection{Kummer valuation tests}

Let \(F=k(C)\) and let
\[
  K=F(y),\qquad y^{31}=x(x-1).
\]

\begin{lemma}[Boundary valuation obstruction]
\label{lem:kummer-valuation}
None of
\[
  r(r-1),\qquad \frac{1-r}{r^2},\qquad \frac{r}{(r-1)^2}
\]
is a \(31\)st power in \(K\).
\end{lemma}

\begin{proof}
Because \(k\) is algebraically closed and contains \(\mu_{31}\), Kummer theory
says that an element \(B\in F^\times\) becomes a \(31\)st power in \(K\) only
if its class in \(F^\times/F^{\times31}\) lies in the subgroup generated by
\(x(x-1)\).  Equivalently,
\[
  B=(x(x-1))^m a^{31}
\]
for some \(m\in\mathbb Z/31\mathbb Z\) and \(a\in F^\times\).  Therefore the
valuation vector of \(B\) on the twelve points of \(U\), modulo \(31\), must
be \(m\) times the valuation vector of \(x(x-1)\), namely
\[
  1\text{ on }E_0,\qquad 1\text{ on }E_1,\qquad -2\text{ on }E_\infty.
\]

The three displayed functions have valuation vectors
\[
\begin{array}{c|ccc}
        &F_0&F_1&F_\infty\\ \hline
r(r-1)&1&1&-2\\
(1-r)/r^2&-2&1&1\\
r/(r-1)^2&1&-2&1 .
\end{array}
\]
If one of these vectors equals \(m(1,1,-2)\) after pulling back to the
\(E\)-partition of \(U\), then the unique four-point block on which the value
is \(-2\) must be \(E_\infty\).  The alternative \(m\ne1\) is impossible:
then the two values \(m\) and \(-2m\) would have to be among \(\{1,-2\}\), and
the only such scalar preserving the multiplicities \(8\) and \(4\) is
\(m=1\).

Thus the three necessary conditions are respectively
\[
  F_\infty=E_\infty,\qquad F_0=E_\infty,\qquad F_1=E_\infty.
\]
For profile \(4\), no column of \(M\) is \((0,0,4)^t\), so none of
\(F_0,F_1,F_\infty\) equals \(E_\infty\).  All three elements therefore fail
the valuation test.
\end{proof}

\subsection{A characteristic-five \(p\)-basis constraint}

Let \(F=k(C)\).  Since \(x\) is separable and \(k\) is perfect of
characteristic \(5\),
\[
  F=F^5(x).
\]
Thus
\[
  r=u_0^5+u_1^5x+u_2^5x^2+u_3^5x^3+u_4^5x^4
\]
for unique \(u_i\in F\), and
\[
  \frac{dr}{dx}=u_1^5+2u_2^5x+3u_3^5x^2+4u_4^5x^3.
\]
At an unramified \(x\)-pole \(P\in E_\infty\), the valuations of the five
terms \(u_i^5x^i\) have distinct residues modulo \(5\).  Therefore there is
no cancellation among them.

As a check on possible characteristic-\(5\) constructions, suppose
\[
  r=x+u^5
\]
and suppose \(r\) is again a profile map with branch values \(0,1,\infty\).
Any pole of \(u\) away from \(x^{-1}(\infty)\) would give \(r\) a pole whose
order is divisible by \(5\), which is incompatible with the profile.  At the
four simple poles in \(E_\infty\), the same residue-modulo-\(5\) argument shows
that \(u\) is regular.  At \(P_\infty\), the pole order of \(u\) is at most
\(6\), otherwise \(u^5\) would dominate the order-\(31\) pole of \(x\).
Thus \(\deg u\le6\).  At a finite high point of \(x\), say
\(x=a+t^{31}\), the equality
\[
  r-r(P)=t^{31}+(u-u(P))^5
\]
has tame ramification index \(31\) only if \(\ord_P(u-u(P))\ge7\), unless
\(u\) is constant.  This contradicts \(\deg u\le6\).  Hence \(u\) is constant,
and then the branch-value condition forces that constant to be \(0\).

\section{High-point coincidences}
\label{sec:highpoint}

A high-point coincidence is an equality \(Q_j=P_i\) among the six ramified
points of the two profile maps.  Its entry is \(M_{ij}\); for instance an
entry-one coincidence means \(M_{ij}=1\).

\begin{lemma}[High-point quotient]
\label{lem:highpoint-quotient}
If \(Q_j=P_i\), then
\[
  Y_{i,j}:=\frac{\Omega_j(r)}{\Omega_i(x)}
\]
has degree at most \(6-M_{ij}\).
\end{lemma}

\begin{proof}
Let \(\{a,b,i\}=\{0,1,\infty\}\) and
\(\{c,d,j\}=\{0,1,\infty\}\).  By Lemma~\ref{lem:omega-divisor}, the order
\(30\) terms at the common point \(Q_j=P_i\) cancel, and
\[
  \Div(Y_{i,j})
  =
  (P_a+P_b+F_j)-(Q_c+Q_d+E_i).
\]
The positive part has degree at most
\[
  2+\deg(F_j-E_i)_+=2+4-M_{ij}=6-M_{ij}.
\]
Any additional high-point coincidence only cancels more terms.
\end{proof}

At every atom \(U_{\alpha,\beta}\) with \(\alpha\ne i\) and \(\beta\ne j\),
both logarithmic forms have simple poles, so the value of \(Y_{i,j}\) is the
quotient of the two residues in the residue table.  Counting these reduced
atoms gives the following forced finite fibers:
\[
\begin{array}{c|c|c|c}
\text{coincidence} & M_{ij} & \deg Y_{i,j}\le & \text{forced finite values}\\ \hline
Q_0=P_0       &0&6& 1:2,\ -1:2\\
Q_1=P_0       &1&5& 1:1,\ -1:4\\
Q_\infty=P_0  &3&3& 1:5,\ -1:2\\
Q_0=P_1       &1&5& 1:1,\ -1:4\\
Q_1=P_1       &2&4& -1:6\\
Q_\infty=P_1  &1&5& 1:4,\ -1:1\\
Q_0=P_\infty  &3&3& 1:5,\ -1:2\\
Q_1=P_\infty  &1&5& 1:4,\ -1:1\\
Q_\infty=P_\infty&0&6&1:2,\ -1:2 .
\end{array}
\]
Hence the coincidences
\[
  Q_\infty=P_0,\qquad Q_1=P_1,\qquad Q_0=P_\infty
\]
are impossible.

\begin{lemma}[Simultaneous inversion]
\label{lem:inversion}
The case \(Q_\infty=P_\infty\) is equivalent to the case \(Q_0=P_0\).
\end{lemma}

\begin{proof}
Replace \(x,r\) by
\[
  x'=1/x,\qquad r'=1/r.
\]
Then
\[
  \Div(x')=31P_\infty+E_\infty-31P_0-E_0,
\]
and
\[
  \Div(x'-1)=\Div((1-x)/x)
  =31P_1+E_1-31P_0-E_0.
\]
Thus the labels \(0\) and \(\infty\) are interchanged for \(x\), while the
label \(1\) is fixed.  The same calculation applies to \(r\).  The incidence
matrix is unchanged by simultaneous interchange of rows \(0,\infty\) and
columns \(0,\infty\):
\[
  \begin{pmatrix}
  0&1&3\\1&2&1\\3&1&0
  \end{pmatrix}
  \longmapsto
  \begin{pmatrix}
  0&1&3\\1&2&1\\3&1&0
  \end{pmatrix}.
\]
Therefore \(Q_\infty=P_\infty\) for \((x,r)\) becomes \(Q'_0=P'_0\) for
\((x',r')\), and the operation is an involution.
\end{proof}

\begin{lemma}[The paired entry-zero edge is impossible]
\label{lem:paired-edge}
There is no profile-\(4\) pair with
\[
  Q_0=P_0,\qquad Q_\infty=P_\infty.
\]
\end{lemma}

\begin{proof}
Assume both equalities.  Put
\[
  u=\frac{\Omega_0(r)}{\Omega_0(x)},\qquad
  v=\frac{\Omega_\infty(r)}{\Omega_\infty(x)}.
\]
The divisor formula gives
\[
  \Div(u)=P_1+U_{1,0}+U_{\infty,0}
          -Q_1-U_{0,1}-U_{0,\infty},
\]
and
\[
  \Div(v)=P_1+U_{0,\infty}+U_{1,\infty}
          -Q_1-U_{\infty,0}-U_{\infty,1}.
\]
Thus both \(u\) and \(v\) have degree \(5\).

Let \(\Gamma\subset\PP^1_u\times\PP^1_v\) be the image of
\((u,v)\), and let \(n=[k(C):k(u,v)]\).  Since both functions have degree
\(5\), the bidegree of \(\Gamma\) is \((5/n,5/n)\), so \(n=1\) or \(5\).
If \(n=5\), then \(\Gamma\) has bidegree \((1,1)\), hence
\(v=M(u)\) for a linear fractional transformation.  The points \(P_1,Q_1\), and the two
points of \(U_{1,1}\), force
\[
  M(0)=0,\qquad M(\infty)=\infty,\qquad M(1)=1,
\]
so \(M\) is the identity.  But \(U_{0,\infty}\) maps to
\((\infty,0)\), contradiction.  Therefore \(n=1\), and \(C\) is the
normalization of a bidegree-\((5,5)\) curve \(\Gamma\).

The atoms \(U_{0,\infty}\), \(U_{\infty,0}\), and \(U_{1,1}\) map to
\[
  (\infty,0),\qquad (0,\infty),\qquad (1,1),
\]
with branch counts \(3,3,2\): the corresponding atoms are reduced on the
normalization \(C\), and \(C\to\Gamma\) is birational.  A singularity of a
reduced curve on a smooth surface with \(m\) distinct analytic branches has
\(\delta\ge {m\choose 2}\).  Hence these three points contribute at least
\[
  {3\choose2}+{3\choose2}+{2\choose2}=7
\]
to the total \(\delta\)-invariant.  The arithmetic genus of a
bidegree-\((5,5)\) curve in \(\PP^1\times\PP^1\) is \(16\), so
\[
  g(C)\le 16-7=9,
\]
contradicting Lemma~\ref{lem:profile-genus}.
\end{proof}

\begin{corollary}[Entry-zero coverage modulo entry-one]
\label{cor:entry0-coverage}
Assume \(Q_0=P_0\).  If no entry-one high-point coincidence occurs, then no
other high-point coincidence occurs.
\end{corollary}

\begin{proof}
The high-point quotient table already excludes
\[
  Q_\infty=P_0,\qquad Q_1=P_1,\qquad Q_0=P_\infty.
\]
The equalities \(Q_1=P_0\) and \(Q_0=P_1\) would identify distinct high points
of a single map.  The remaining possible extra coincidences are
\[
  Q_1=P_\infty,\qquad Q_\infty=P_1,\qquad Q_\infty=P_\infty.
\]
The first two are entry-one coincidences, and the last is impossible by
Lemma~\ref{lem:paired-edge}.
\end{proof}

\section{The entry-one branch}
\label{sec:entry-one}

The four entry-one coincidences are equivalent under the symmetries of the
incidence matrix and by interchanging the two maps.  We use the representative
\[
  Q_1=P_0.
\]

\subsection{The degree-five quotient}

Define
\[
  z=\frac{\Omega_1(r)}{\Omega_0(x)}
   =\frac{dr/r}{dx/(x-1)}.
\]
By Lemma~\ref{lem:omega-divisor},
\[
  \Div(z)=P_1+P_\infty+D+G-Q_0-Q_\infty-B.
\]
Thus \(\deg z=5\) in the no-further-highpoint subcase.  If another high-point
coincidence is also present, the divisor above has further cancellation; such
lower-degree specializations remain part of the entry-one problem.  We first
treat the no-further-highpoint case.

The boundary residues force
\[
  z=1\text{ on }C,\qquad z=-1\text{ on }E+F.
\]
Let
\[
  \alpha=z(P_0)=z(Q_1),\qquad u=z(A),\qquad
  c=x(Q_0),\qquad d=x(Q_\infty).
\]
The points \(Q_0\) and \(Q_\infty\) may lie in the same \(x\)-fiber; no
condition \(c\ne d\) is imposed.

Since \(\deg x=35\) and \(\deg z=5\), if \(k(x,z)\ne k(C)\) then the degree
\([k(C):k(x,z)]\) divides both \(35\) and \(5\).  The nontrivial possibility
would force \(x\in k(z)\), so the ramification index \(31\) of \(x\) would
factor through a degree-\(5\) map, impossible.  Hence
\[
  k(C)=k(x,z),
\]
and \(C\) is the normalization of an integral curve of bidegree \((5,35)\)
in \(\PP^1_x\times\PP^1_z\).

\subsection{Norms and the tautological trace identity}

Write \(\Norm_z=\Norm_{k(C)/k(z)}\).  Pushing the divisors of
\(x,x-1,r,r-1\) along \(z\) gives, up to constants,
\[
  \Norm_z(x)=\frac{(z-\alpha)^{31}(z-u)}{z^{32}(z+1)^3},
\qquad
  \Norm_z(x-1)=\frac{z(z-1)}{(z+1)^2},
\]
and
\[
  \Norm_z(r)=(z-1)(z+1)^2,\qquad
  \Norm_z(r-1)=\frac{(z-\alpha)^{31}(z-u)z^3}{z+1}.
\]
The defining differential relation is
\[
  d\log r=z\,d\log(x-1).
\]
Taking trace to \(k(z)\) gives
\[
  d\log\Norm_z(r)=z\,d\log\Norm_z(x-1).
\]
Substitution yields
\[
  \frac1{z-1}+\frac2{z+1}
  =
  z\left(\frac1z+\frac1{z-1}-\frac2{z+1}\right),
\]
which is an identity.  Thus the first pushed-forward norm/trace relation is
not an obstruction.

\subsection{The \(p\)-closed equation}

Let
\[
  \partial=(x-1)\frac{d}{dx}.
\]
For this entry-one quotient, the Cartier-operator formula for the logarithmic
exact differential \(d\log r=z\,d\log(x-1)\) gives the characteristic-\(5\)
equation
\[
  \partial^4z=z-z^5.
\]
This at least proves separability of \(z\).  If \(dz=0\), then
\(\partial z=0\), so the equation gives \(z^5=z\).  Over \(k=\kbar\), this
forces \(z\in\F_5\), contradicting \(\deg z=5\).

\subsection{A bidegree-\((5,35)\) normal form}

Normalize the bidegree equation \(f(x,z)=0\) by
\[
  [z^{35}]f=x^3(x-c)(x-d),\qquad
  [x^5]f=z^{32}(z+1)^3.
\]
Put
\[
  s=c+d,\qquad p=cd,\qquad N=(1-c)(1-d).
\]
Here \(N\ne0\); otherwise the fiber \(x=1\) would contain a pole of \(z\),
contradicting the displayed divisor of \(z\).
Let
\[
  \Phi(z)=(z-\alpha)^{31}(z-u),\quad
  \Psi(z)=z^{33}(z-1)(z+1),\quad
  \Gamma(z)=z^{32}(z+1)^3.
\]
Then one may write
\[
  f(x,z)=(1-x)K\Phi(z)+xN\Psi(z)+x(x-1)H(x,z),
\]
where
\[
  H=h_0(z)+xh_1(z)+x^2h_2(z)+x^3\Gamma(z).
\]
The leading \(z^{35}\)-coefficients are forced:
\[
  [z^{35}]h_0=N,\qquad [z^{35}]h_1=N,\qquad
  [z^{35}]h_2=1-s.
\]
Indeed the coefficient of \(z^{35}\) is
\[
  xN+x(x-1)\bigl(N+Nx+(1-s)x^2+x^3\bigr)
  =x^3(x^2-sx+p).
\]

The fiber \(x=0\) gives
\[
  f(0,z)=K(z-\alpha)^{31}(z-u),
\]
the fiber \(x=1\) gives
\[
  f(1,z)=N z^{33}(z-1)(z+1),
\]
and the \(x=\infty\) boundary is encoded by \([x^5]f=\Gamma(z)\).

At the corner \((x,z)=(0,\infty)\), write \(w=1/z\) and
\[
  F(x,w)=w^{35}f(x,1/w).
\]
The three points of \(B=U_{0,\infty}\) have local behavior \(w\sim x\).  Thus
the tangent cubic is proportional to
\[
  (w-x)^3.
\]
Writing
\[
  w^{35}h_0(1/w)=N+aw+bw^2+O(w^3),\qquad
  w^{35}h_1(1/w)=N+ew+O(w^2),
\]
the degree-two term forces \(a=0\), and the cubic comparison gives
\[
  p=-K,\qquad e=-3K,\qquad b=3K-N.
\]
These are genuine constraints, but they are not yet contradictory.

\begin{opentheorem}[B. Entry-one exclusion]
\label{open:entry-one}
No profile-\(4\) pair has an entry-one high-point coincidence.  In the
representative \(Q_1=P_0\) with no further high-point coincidence, this is the
assertion that no integral bidegree-\((5,35)\) curve \(f(x,z)=0\), with the
boundary normal form and tangent constraints above, is the normalization of a
profile-\(4\) pair satisfying
\[
  d\log r=z\,d\log(x-1)
\]
for some function \(r\) with the prescribed divisor.  The assertion includes
the lower-degree boundary strata obtained from this representative by imposing
additional high-point coincidences compatible with the quotient table.
\end{opentheorem}

\section{The no-highpoint branch}
\label{sec:no-highpoint}

If none of the six high points \(P_0,P_1,P_\infty,Q_0,Q_1,Q_\infty\) agree,
the low-degree quotient of Lemma~\ref{lem:highpoint-quotient} is unavailable.
The residual approach below uses a bidegree-\((4,4)\) ansatz.  The first
unproved part is that every no-highpoint solution is represented by this
ansatz.

Let
\[
\begin{aligned}
P(X,R)=&\ R^4X^4-R^4X^3-R^3X^4
       +2R^3X-2R^2X\\
      &\quad +2RX^3-2RX^2+R+X-1,
\end{aligned}
\]
and
\[
  L=4XR+4X+4R+2.
\]
The line \(L=0\) is the graph
\[
  R=\frac{4X+2}{X+1},
\]
which sends \(X=0,1,\infty\) to \(R=2,3,4\), respectively.  The boundary
preserving residual family is
\[
  H_A=L^{31}P-X(X-1)R(R-1)A(X,R),
\]
where
\[
  \deg_X A\le 32,\qquad \deg_R A\le 32.
\]
The factor \(X(X-1)R(R-1)\) makes the six affine boundary restrictions agree
with those of \(L^{31}P\).  The condition that the
\((X,R)=(0,\infty)\) and \((\infty,0)\) corners have multiplicity three
forces the two corner coefficients
\[
  [R^{32}]A(0,R)=2,\qquad [X^{32}]A(X,0)=2.
\]

Given the boundary profile, Riemann--Hurwitz says that a valid residual curve
has no interior ramification for either projection: the three order-\(31\)
boundary branches already use the full ramification budget.  Thus a residual
candidate must satisfy, away from the boundary,
\[
  H_A=\frac{\partial H_A}{\partial R}=0
  \quad\text{has no solution}
\]
for the \(X\)-projection, and similarly with \(X\) and \(R\) interchanged.
In characteristic \(5\), writing \(D=X(X-1)R(R-1)\) and \(F=LP\), the
ramification equations away from \(F=0\) may be written as
\[
  N_R=F\,\partial_R(DA)-DA\,\partial_RF=0,
\]
\[
  N_X=F\,\partial_X(DA)-DA\,\partial_XF=0.
\]

Several small subfamilies of this residual system have been eliminated:
constant \(A\), sparse low-parameter families, and the symmetric
ordinary-corner candidates.  These eliminations are not enough for the full
coefficient space, which has \(33^2\) coefficients before boundary conditions,
and they are not used as proof of the full no-highpoint branch below.

\begin{opentheorem}[C. No-highpoint exclusion]
\label{open:no-highpoint}
No profile-\(4\) two-map pair exists with all six high points
\[
  P_0,P_1,P_\infty,Q_0,Q_1,Q_\infty
\]
distinct.

A sufficient residual route to this theorem is the conjunction of the
following two precise statements:
\begin{enumerate}[label=\textnormal{(\arabic*)},leftmargin=2.2em]
\item every no-highpoint profile-\(4\) pair gives a residual equation
  \(H_A=0\) of the form above;
\item no polynomial \(A\) in that residual family satisfies all boundary,
  integrality, genus, and no-interior-ramification conditions over \(\kbar\).
\end{enumerate}
\end{opentheorem}

\section{The entry-zero branch}
\label{sec:entry-zero}

By Lemma~\ref{lem:inversion}, it is enough to study the entry-zero
representative
\[
  Q_0=P_0.
\]
By Corollary~\ref{cor:entry0-coverage}, after the entry-one branch is removed,
this representative has no other high-point coincidence.

\subsection{The quotient and its norms}

Assume \(Q_0=P_0\) and no other high-point coincidence.  Define
\[
  z=\frac{\Omega_0(r)}{\Omega_0(x)}
   =\frac{dr/(r-1)}{dx/(x-1)}.
\]
Then
\[
  \Div(z)=P_1+P_\infty+C+F-Q_1-Q_\infty-A-B,
\]
so \(\deg z=6\).  Since \(\gcd(35,6)=1\), the field \(k(x,z)\) is all of
\(k(C)\).  Therefore \(C\) is the normalization of an integral bidegree
\((6,35)\) curve \(f(x,z)=0\) in \(\PP^1_x\times\PP^1_z\).

Let
\[
  \alpha=z(P_0),\qquad c=x(Q_1),\qquad d=x(Q_\infty),
\]
and set
\[
  s=c+d,\qquad p=cd,\qquad N=(1-c)(1-d)=1-s+p.
\]
The no-other-highpoint assumption and the fact that \(U\) is the common
unramified boundary imply
\[
  c,d\in k\setminus\{0,1\};
\]
the points \(Q_1,Q_\infty\) are distinct, but the values \(c,d\) need not be
distinct.
Normalize \(f\) by
\[
  [x^6]f=z^{34}(z+1),\qquad [z^{35}]f=x^4(x-c)(x-d).
\]
The boundary fibers have the form
\[
  f(x,0)=K(x-1)^2,
\]
\[
  f(0,z)=K(1-\eta z)^{31},\qquad \eta=\alpha^{-1},
\]
and
\[
  f(1,z)=N z^{32}(z-1)^2(z+1).
\]
These formulas are just the pushed-forward divisors of \(z,x,x-1\).  Explicitly,
\[
  \Norm_{C/k(x)}(z)=
  \frac{(x-1)^2}{x^4(x-c)(x-d)},
\]
\[
  \Norm_{C/k(z)}(x)=
  \frac{(z-\alpha)^{31}}{z^{34}(z+1)},
\]
and
\[
  \Norm_{C/k(z)}(x-1)=\frac{(z-1)^2}{z^2}.
\]

\begin{lemma}[Norm-constant reduction]
\label{lem:cd2}
In the entry-zero no-other-highpoint branch,
\[
  \alpha^{31}=1,\qquad (1-c)(1-d)=1,\qquad cd=-1.
\]
Consequently
\[
  c=d=2.
\]
\end{lemma}

\begin{proof}
Since \(f\) has odd degree \(35\) in \(z\),
\[
  \Norm_{C/k(x)}(z)=
  -\frac{f(x,0)}{x^4(x-c)(x-d)}
  =
  -\frac{K(x-1)^2}{x^4(x-c)(x-d)}.
\]
Comparison with the exact norm gives \(K=-1\).

Since \(f\) has degree \(6\) in \(x\),
\[
  \Norm_{C/k(z)}(x)=
  \frac{f(0,z)}{z^{34}(z+1)}
  =
  \frac{K(1-\eta z)^{31}}{z^{34}(z+1)}.
\]
As \(1-\eta z=-\eta(z-\alpha)\), comparison gives
\[
  K(-\eta)^{31}=1.
\]
With \(K=-1\), this implies \(\eta^{31}=1\), hence \(\alpha^{31}=1\).

Similarly,
\[
  \Norm_{C/k(z)}(x-1)=
  \frac{f(1,z)}{z^{34}(z+1)}
  =
  N\frac{(z-1)^2}{z^2},
\]
so \(N=1\).

It remains to identify \(p=cd\).  Put \(w=1/z\) and
\[
  \bar f(x,w)=w^{35}f(x,1/w).
\]
At \((x,w)=(0,0)\), the residues of
\((dr/(r-1))/(dx/(x-1))\) give \(z\sim -1/x\) on the atom
\(A=U_{0,1}\), while the three points of \(B=U_{0,\infty}\) have
\(z\sim 1/x\).  Hence the degree-four
tangent cone is proportional to
\[
  (w+x)(w-x)^3.
\]
The \(w^4\)-coefficient of the degree-four part of \(\bar f\) is
\(K(-\eta)^{31}\), and the \(x^4\)-coefficient is \(p\).  In
\((w+x)(w-x)^3\), these two coefficients have opposite signs.  Therefore
\[
  p=-K(-\eta)^{31}=-1.
\]
Since \(N=1-s+p=1\) and \(p=-1\), we have \(s=-1\).  The numbers \(c,d\) are
the roots of
\[
  T^2-sT+p=T^2+T-1.
\]
Over characteristic \(5\), this is \((T-2)^2\).  Hence \(c=d=2\).
\end{proof}

Thus the entry-zero branch is reduced to the corrected double-fiber case
\[
  Q_0=P_0,\qquad x(Q_1)=x(Q_\infty)=2,\qquad Q_1\ne Q_\infty.
\]

\subsection{The formal local tower}

The double-fiber case is studied by expanding the normalized bidegree
\((6,35)\) equation at the double fiber \(x=2\) and imposing the
characteristic-\(5\) differential equation attached to
\[
  z=\frac{dr/(r-1)}{dx/(x-1)}.
\]
The expansion alternates repeated and simple layers.  The layers through
simple \(u^{25}\), together with the repeated \(u^{25}\) two-minor step, have
explicit linear certificates: at each layer the relevant residual equations
either solve the next group of coefficients or cut to a smaller affine
subspace by a displayed full-rank matrix.  The first unresolved global local
frontier is the tail-exhausted simple-\(u^{30}\) all-zero locus together with
the base residual and repeated-\(e30\) consistency.

One checked component of this frontier is the point \(P129\).  In local
coordinates, the representative is
\[
(2,1,2,3,4,3,2,1,1,3,0,2,0,2,4,2,1,1,2).
\]
After the repeated-\(u^{25}\) pivot solve, the remaining free variables are
\[
\begin{gathered}
u7\_3,u6\_2,u6\_3,u5\_1,u5\_2,u4\_0,u4\_1,u4\_2,u4\_3,\\
u3\_0,u3\_1,u3\_2,u3\_3,u2\_2,u1\_0,u1\_1 .
\end{gathered}
\]
The repeated-\(e30\) condition is the affine three-row system
\[
  (3,2,4)+\sum_{\nu=1}^{16}y_\nu c_\nu=0,
\]
where the columns \(c_\nu\), in the order above, are
\[
\begin{array}{c|c}
\text{variable}&c_\nu\\ \hline
u7\_3&(1,4,1)\\
u6\_2&(4,3,3)\\
u6\_3&(0,1,1)\\
u5\_1&(2,2,4)\\
u5\_2&(0,1,3)\\
u4\_0&(2,3,3)\\
u4\_1&(2,1,2)\\
u4\_2&(0,1,4)\\
u4\_3&(3,1,1)\\
u3\_0&(1,4,4)\\
u3\_1&(1,1,1)\\
u3\_2&(4,3,0)\\
u3\_3&(1,4,3)\\
u2\_2&(0,3,2)\\
u1\_0&(1,0,0)\\
u1\_1&(3,3,3).
\end{array}
\]
This matrix has rank \(3\), leaving an affine \(13\)-space.  The later
layers \(e35,e40,e45\) have ranks
\[
  4,\qquad 4,\qquad 4
\]
on successive solution spaces, leaving a final line.  On that line the simple
\(e50\) residual is the nonzero constant \(2\).  Therefore \(P129\) does not
extend to a formal local solution through \(e50\).  The same \(e50\) ending
also kills the checked nearby tangent points of \(P129\) and the checked
\(\rho4\) tangent hit.

\subsection{Basin \(127\)}

Another high-nullity component of the same frontier is a six-parameter affine
family with coordinates \(r_0,\ldots,r_5\).  The expected saturated quotient is
the ideal
\[
\begin{gathered}
J=\langle
r_2-2,\ r_5-2,\ 
r_0+2r_1-2r_3+2r_4^2+r_4+1,\\
r_1^2-r_3-2r_4-1,\ 
r_1r_3+2r_1-2r_3-r_4^2+2r_4,\\
r_1r_4-r_1+2r_4^2-2,\ 
r_3^2+r_3-2,\ 
r_3r_4-r_3-r_4+1,\\
r_4^3-2r_4^2-r_4+2
\rangle
\subset \F_5[r_0,\ldots,r_5].
\end{gathered}
\]
\begin{lemma}[The displayed basin-\(127\) quotient]
\label{lem:J127-points}
The ideal \(J\) is radical of length \(5\).  Its five points are
\[
\begin{gathered}
(0,0,2,1,4,2),\quad
(0,1,2,3,1,2),\quad
(3,4,2,1,2,2),\\
(4,2,2,1,1,2),\quad
(4,4,2,3,1,2).
\end{gathered}
\]
Moreover, adding the base residual
\[
  B_{127}=r_3+r_4^2.
\]
cuts \(J\) to the single maximal ideal
\[
  (r_0+2,\ r_1+1,\ r_2-2,\ r_3-1,\ r_4-2,\ r_5-2),
\]
which is the already killed \(P129\) point.
\end{lemma}

\begin{proof}
The generators first give \(r_2=r_5=2\) and
\[
  r_4^3-2r_4^2-r_4+2=(r_4-1)(r_4-2)(r_4+1).
\]
Also
\[
  r_3r_4-r_3-r_4+1=(r_4-1)(r_3-1).
\]
If \(r_4=4=-1\), then \(r_3=1\).  The equation
\[
  r_1r_4-r_1+2r_4^2-2=0
\]
gives \(r_1=0\), and the linear equation for \(r_0\) gives \(r_0=0\).  If
\(r_4=2\), then \(r_3=1\), the same equation gives \(r_1=4\), and then
\(r_0=3\).

It remains to set \(r_4=1\).  Then \(r_3^2+r_3-2=0\), so
\[
  r_3\in\{1,3\}.
\]
For \(r_3=1\), the equation
\[
  r_1r_3+2r_1-2r_3-r_4^2+2r_4=0
\]
gives \(r_1=2\), and then \(r_0=4\).  For \(r_3=3\), the equation
\[
  r_1^2-r_3-2r_4-1=0
\]
gives \(r_1\in\{1,4\}\), and the linear equation gives respectively
\(r_0=0\) and \(r_0=4\).  These are exactly the five displayed points.

All factors in this triangular decomposition are squarefree over \(\F_5\), so
\(J\) is the intersection of the corresponding maximal ideals.  Evaluating
\(B_{127}=r_3+r_4^2\) on the five points leaves only
\((3,4,2,1,2,2)\), which is the maximal ideal displayed above.
\end{proof}

\begin{target}[D1. Basin-\(127\) containment]
\label{target:basin127}
Let \(I_{127}\) be the ideal generated by the basin-\(127\) low-data equations
in the six coordinates \(r_0,\ldots,r_5\), and let \(s_{127}\) be the product
of the etale discriminant and the chart pivot.  Then
\[
  V(I_{127})\cap D(s_{127})\subseteq V(J),
\]
or equivalently every generator of \(\sqrt{J}\) vanishes on the localized
low-data scheme.  The stronger saturated equality
\[
  I_{127}:s_{127}^{\infty}=J
\]
implies this containment.
\end{target}

Target~\ref{target:basin127} is the part presently being checked in Sage.  It
would close the basin-\(127\) component by reducing its base-residual locus to
the already killed \(P129\) point.

\begin{target}[D2. Double-fiber tail coverage]
\label{target:tail-coverage}
On each \(\rho3\) and \(\rho4\) pivot chart, form the localized ideal
\[
  I_{\mathrm{chart}}
  =
  \langle
  \Phi_0,\Phi_1,\Phi_2,\,
  S_{30},\,
  E30_0,E30_1,E30_2
  \rangle:s_{\mathrm{chart}}^\infty,
\]
where \(\Phi_i\) are the shared repeated low-data equations, \(S_{30}\) is the
simple-\(u30\) Schur/base residual, \(E30_i\) are the three affine repeated
\(e30\) equations after the repeated-\(u25\) pivot solve, and
\(s_{\mathrm{chart}}\) is the product of the etale discriminant and all chart
pivots.  Then \(V(I_{\mathrm{chart}})\) is contained in the union of
the \(P129\), \(\rho4\), and basin-\(127\) strata, and after the
corresponding \(e35,e40,e45,e50\) layer equations are adjoined on each stratum,
the remaining localized ideals are unit ideals.
\end{target}

\begin{opentheorem}[D. Double-fiber completion]
\label{open:double-fiber}
There is no profile-\(4\) pair satisfying
\[
  Q_0=P_0,\qquad
  x(Q_1)=x(Q_\infty)=2,\qquad
  Q_1\ne Q_\infty,
\]
with no other high-point coincidence.  It follows from the two certificate
targets above.
\end{opentheorem}

\section{Assembly of the curve-level reduction}
\label{sec:assembly}

\begin{proposition}[Certificate targets imply Open Theorem D]
\label{prop:entry0-conditional}
Assume Certificate Targets~\ref{target:basin127} and
\ref{target:tail-coverage}.  Then Open Theorem~\ref{open:double-fiber}
holds.
\end{proposition}

\begin{proof}
By Lemma~\ref{lem:cd2}, every entry-zero no-other-highpoint solution lies in
the corrected \(c=d=2\) double-fiber branch.  The formal layer certificates
through repeated \(u^{25}\) and simple \(u^{30}\) reduce this branch to the
tail-exhausted simple-\(u30\) all-zero locus with base residual and
repeated-\(e30\) consistency.  The \(P129\) and \(\rho4\) local strata are
killed at \(e50\).  Target~\ref{target:basin127} reduces the basin-\(127\)
base-residual part to \(P129\), already killed.  Target~\ref{target:tail-coverage}
says that every other point of the localized frontier is contained in one of
these killed strata or becomes empty after the later layer equations.
\end{proof}

\begin{proposition}[Open Theorems B--D eliminate profile \(4\)]
\label{prop:BCD-profile4}
Assume Open Theorems~\ref{open:entry-one}, \ref{open:no-highpoint}, and
\ref{open:double-fiber}.  Then no profile-\(4\) two-map pair exists.
\end{proposition}

\begin{proof}
Let \((C,x,r)\) be a profile-\(4\) pair.  If there is no high-point
coincidence, Open Theorem~\ref{open:no-highpoint} removes it.  If there is an
entry-one high-point coincidence, Open Theorem~\ref{open:entry-one} removes it.  The
high-point quotient table removes the three coincidences
\[
  Q_\infty=P_0,\qquad Q_1=P_1,\qquad Q_0=P_\infty.
\]
Thus any remaining high-point coincidence is entry-zero.  By simultaneous
inversion it is represented by \(Q_0=P_0\).  After entry-one cases have been
removed, Corollary~\ref{cor:entry0-coverage} says there is no other high-point
coincidence.  Lemma~\ref{lem:cd2} reduces this branch to the corrected
double-fiber case, and Open Theorem~\ref{open:double-fiber} removes it.
Therefore no profile-\(4\) pair remains.
\end{proof}

\begin{theorem}[Visibility from the four open theorems]
\label{thm:four-open-imply-comm}
Assume Open Theorems~\ref{open:profile-extraction}, \ref{open:entry-one},
\ref{open:no-highpoint}, and \ref{open:double-fiber}.  Then every finite
etale self-correspondence of \(S=\PP^1_k(31,31,31)\) is visible.  Equivalently,
\(\Comm_{\mathrm{alg},k}(S)=\Aut(S)=S_3\).
\end{theorem}

\begin{proof}
Let \(u,v:T\to S\) be a connected representable finite etale
self-correspondence.  If
\[
  \pi_0\circ u\not\simeq \pi_0\circ v,
\]
then Open Theorem~\ref{open:profile-extraction} produces a profile-\(4\)
two-map pair.  This contradicts Proposition~\ref{prop:BCD-profile4}.  Hence
\(\pi_0\circ u\simeq\pi_0\circ v\) for every connected finite etale
self-correspondence of \(S\), so every such correspondence is visible by
Definition~\ref{def:visible}.
\end{proof}

\begin{corollary}[Conditional common-cover counterexample from the four open theorems]
If Open Theorems~\ref{open:profile-extraction}, \ref{open:entry-one},
\ref{open:no-highpoint}, and \ref{open:double-fiber} are proved, then
\[
  Y:\ y^{31}=x(x-1)
\]
has no connected common finite etale cover with any smooth proper connected
genus-three curve over \(k\).
\end{corollary}

\begin{proof}
Combine Theorem~\ref{thm:four-open-imply-comm} with
Proposition~\ref{prop:conditional-common-cover}.
\end{proof}

\section{Open Theorems}

The reduction leaves the following four independent statements.
\begin{enumerate}[label=\textnormal{(\Alph*)},leftmargin=2.2em]
\item Open Theorem~\ref{open:profile-extraction}: every non-visible
  finite etale self-correspondence of \(S\) produces the profile-\(4\)
  degree-\(35\) pair.
\item Open Theorem~\ref{open:entry-one}: no profile-\(4\) pair has an
  entry-one high-point coincidence.
\item Open Theorem~\ref{open:no-highpoint}: no profile-\(4\) pair has
  all six high points distinct.
\item Open Theorem~\ref{open:double-fiber}: the corrected entry-zero
  double-fiber branch is empty.
\end{enumerate}

The Sage computation concerns item (D), more specifically Certificate
Target~\ref{target:basin127}.  If it proves
\[
  V(I_{127})\cap D(s_{127})\subseteq V(J)
\]
or the stronger saturation equality
\[
  I_{127}:s_{127}^{\infty}=J,
\]
then the basin-\(127\) component of the double-fiber frontier is closed: the
base residual cuts it to \(P129\), already killed at \(e50\).  That still
leaves Certificate Target~\ref{target:tail-coverage}, the chart-wise coverage
of the rest of the tail-exhausted simple-\(u30\) base-zero/e30 locus.

Thus the computation is one algebraic containment inside Open
Theorem~\ref{open:double-fiber}.  Together with Open
Theorems~\ref{open:profile-extraction}, \ref{open:entry-one}, and
\ref{open:no-highpoint}, Open Theorem~\ref{open:double-fiber} implies the
visibility theorem by Theorem~\ref{thm:four-open-imply-comm}.

\end{document}
